\documentclass[exercice]{thedocument}%
\startdatakeys
\source{}
\cycle{}
\competencesDuSocle{}
\connaissancesRequises{}
\competencesTravaillees{}
\enddatakeys
\begin{document}%
Une collectionneuse compte ses cartes Pokémon afin de les revendre. Elle
possède $252$ cartes de type "feu" et $156$ cartes de type "sol".
\begin{enumerate}
    \item
    \begin{enumerate} 
        \item Parmi les trois propositions suivantes, laquelle correspond à la décomposition en produit de facteurs premiers du nombre $252$~:\par 
        \begin{minipage}[t]{.3\linewidth}
            $\square$\, $2^2\times 9\times 7$
        \end{minipage}\hfill
        \begin{minipage}[t]{.3\linewidth}
            $\square$ \, $2\times 2\times 3\times 21$
        \end{minipage}\hfill
        \begin{minipage}[t]{.3\linewidth}
            $\square$\, $2^2\times 3^2\times 7$
        \end{minipage}
        \item Donner la décomposition en produit de facteurs premiers du nombre $156$.
    \end{enumerate}
        \item Elle veut réaliser des paquets identiques, c'est-à-dire contenant chacun le même nombre de cartes "sol" et le même nombre de cartes "feu" en utilisant toutes ses cartes. 
    \begin{enumerate}
        \item Peut-elle faire $36$ paquets~?
        \item Quel est le nombre maximum de paquets qu'elle peut réaliser~?
        \item Combien de cartes de chaque type contient alors chaque paquet~? 
    \end{enumerate}
\end{enumerate}
\hfill\scriptsize\textit{Source~: Brevet Metropole, juin 2022}
\end{document}
